\chapter*{\chapterstyle{III --- Corps des Réels}} % 95% Fini
\addcontentsline{toc}{section}{Corps des Réels}

Soit \(A\) une partie de \(\R\) et \(M \in \R\), on dit que \(M\) est un \textbf{majorant} de \(A\) si il est est supérieur à tout élément de \(A\), 
en outre, si \(M \in A\), on dit que \(M\) est le \textbf{maximum} de \(A\)\footnote[1]{Les propriétés sont analogues pour les \textbf{minorants} de \(\R\)}. 

\subsection*{\subsecstyle{Borne supérieure {:}}}
Soit \(M^+\) \textbf{l'ensemble des majorants} de la partie \(A\), alors par construction, cet ensemble admet un minimum\footnote[2]{Voir la construction explicite de \(\R\) par les suites de Cauchy, ou par les coupures de Dedekind.}, qu'on appelle \textbf{borne supérieure}, ie on définit:
\customBox{width=4cm}{
    \(\sup(A) := \min(M^+)\)
}

Soit \(A, B\) deux parties bornées de \(\R\), alors l'application qui associe sa borne supérieure à une partie est \textbf{croissante} et celle qui associe sa borne inférieure à une partie est \textbf{décroissante}.
\subsection*{\subsecstyle{Caractérisations {:}}}
Soit \(M\) un majorant de \(A\), alors on peut alors montrer \textbf{la caractérisation séquentielle} de la borne supérieure:
\customBox{width=5cm}{
    \(\exists (u_n) \in A^\N \; ; \; u_n \longrightarrow M\)
}
\begin{center}
    \textit{
        Un majorant \(M\) est la borne supérieure de \(A\), si et seulement si il existe une suite à valeurs dans \(A\) qui tends vers \(M\).
    }
\end{center}
Soit \(\epsilon > 0\), alors on aussi \textbf{une caractérisation analytique} de la borne supérieure:  
\customBox{width=5.5cm}{
    \(\exists a \in A \; ; \; M - \epsilon \leq a \leq M\)
}

\begin{center}
    \textit{
        Un majorant \(M\) est la borne supérieure de \(A\), si tout nombre plus petit n'est plus un majorant.
    }
\end{center}

\subsection*{\subsecstyle{Propriété d'Archimède}}
Soit \(K\) un corps totalement ordonné, alors \(K\) est dit \textbf{archimédien}, si et seulement si on a\footnote[3]{Le cas trivial \(x < 0\) est exclu pour plus de lisibilité.}:
\customBox{width=6.5cm}{
    \(\forall (a, b) \in K_+^* \; , \; \exists n \in \N \; ; \; na > b\)
}

\begin{center}
    \textit{
        Donc \(K\) est archimédien, si tout nombre de \(K\) peut être dépassé par un multiple entier\+
        d'un autre nombre de \(K\).
    }
\end{center}

Une autre interprétation est \textbf{qu'il n'existe pas de nombre infiniment grand ou petit} dans \(\R\), pour le voir il suffit de poser \(a = 1\) ou \(b = 1\).
\pagebreak

\subsection*{\subsecstyle{Intervalles {:}}}
Soit \(a, b \in \R\) tels que \(a < b\), on définit \textbf{l'intervalle fermé} \([a, b]\) l'ensemble:
\[
    [a, b] := \Bigl\{ x\in \R \; ; \; a \leq x \leq b \Bigl\}
\]
De manière analogue et selon les relations d'ordre strictes ou non, on définit les intervalles \textbf{ouverts et semi-ouverts}.\<

On dit plus généralement que les intervalles sont les parties \textbf{convexes} de \(\R\).

\subsection*{\subsecstyle{Valeur absolue {:}}}
On introduit l'application \textbf{valeur absolue} qui sera très utile par la suite pour mesurer des écarts, elle est définie comme suit:
\[
    |x| := \max(x, -x)
\]
Cette application est une \textbf{métrique}, elle est \textbf{multiplicative} et vérifie les \textbf{inégalités triangulaires}:
\begin{multicols}{2}
    \begin{itemize}
        \item \(|x + y| \leq |x| + |y|\)
        \item \(||x| - |y|| \leq |x - y|\)
    \end{itemize}
\end{multicols}

\subsection*{\subsecstyle{Partie entière {:}}}
Soit \(n \in \N\), on appelle \textbf{partie entière} d'un réel \(x\), et on note \(\lfloor x \rfloor\) \textbf{l'unique} entier relatif tel que:
\[
    \lfloor x \rfloor \leq x < \lfloor x \rfloor + 1
\]
La partie entière vérifie la propriété fondamentale suivante:
\customBox{width=3.5cm}{
    \(\lfloor x + n \rfloor = \lfloor x \rfloor + n\)
}
L'existence de la partie entière est garantie par \textbf{la propriété d'Archimède}.

\subsection*{\subsecstyle{Approximation décimale {:}}}
Soit \(n \in \N\), il peut être utile d'approximer un réel à \(n\) chiffres après la virgule, ie à \(10^{-n}\) près, la partie entière nous permet alors de le faire, et on l'obtient en exécutant l'algorithme suivant:
\[
    x \longrightarrow 10^n x \longrightarrow \lfloor 10^n x \rfloor \longrightarrow \frac{\lfloor 10^n x \rfloor}{10^n}
\]
On obtient alors une approximation à \textbf{n chiffre après la virgule} de \(x\) qu'on note \(d_n(x)\) et on a:
\customBox{width=3.5cm}{
    \[d_n(x) = \frac{\lfloor 10^n x \rfloor}{10^n}\]
}

\chapter*{\chapterstyle{III --- Suites numériques}} % 95% Fini
\addcontentsline{toc}{section}{Suites numériques}

On appelle \textbf{suite} une \textbf{famille d'éléments} indexée par une partie \(I\) de \(\N\), et on note une telle famille \((u_n)_{n \in I}\), qu'on abrège souvent en \((u_n)\).\<

Lorsque tout les éléments de la famille appartiennent au \textbf{même ensemble} \(E\), on peut alors assimiler la suite \((u_n)\) à une application de \(\N\) vers \(E\).\<

Soit \((u_n)\), \((v_n)\) et \((w_n)\) trois suites à valeurs dans K.

\subsection*{\subsecstyle{Limite {:}}}
On dit qu'une suite \((u_n)\) tends vers une limite\footnote[1]{Elle peut aussi tendre vers \(\pm \infty\), ce qui équivaut à montrer que \((u_n)\) n'est pas majorée/minorée.} \(l \in K\) si elle vérifie:
\customBox{width=9cm}{
    \(\exists l \in K \; , \; \forall \epsilon \in \R_+ \; , \; \exists N \in \N \; , \; \forall n > N \; ; \; |u_n - l| < \epsilon \)
}

On dit alors qu'elle converge et sa limite est \textbf{unique}, sinon on dit qu'elle diverge.\+
On alors les propriétés fondamentales suivantes:
\begin{align*}
    &\bullet \text{Si \((u_n)\) converge, alors elle est bornée}\\
    &\bullet \text{Si \((u_n) \rightarrow l\), alors \((|u_n|) \rightarrow |l|\)}
\end{align*}

\subsection*{\subsecstyle{Opérations sur les limites {:}}}
On considère que \((u_n)\) et \((v_n)\) sont convergentes de limites finies \(l, l' \in K^*\).\+
Alors le passage à la limite est une opération \textbf{linéaire et multiplicative}, ie elle se comporte conformément à l'intuition.\<

Dans le cas où les limites sont infinies ou nulles, des \textbf{formes indéterminées} peuvent apparaître et nécessitent les précautions usuelles pour conclure.

\subsection*{\subsecstyle{Suites réelles {:}}}
On considère \((u_n)\), \((v_n)\) et \((w_n)\) des suites à valeurs dans \(\R\). La grande spécificité de \(\R\) est de pouvoir parler de suites \textbf{monotones} (croissantes ou décroissantes) et de suites \textbf{majorées ou minorées}.\<

On a alors le \textbf{théorème de la limite monotone}:
\customBox{width=15cm}{
   Si la suite \((u_n)\) est \textbf{croissante et majorée} alors elle converge vers \(\sup\{ u_n \; ; \; n \in \N\}\).\\
   Si la suite \((u_n)\) est \textbf{décroissante et minorée} alors elle converge vers \(\inf\{ u_n \; ; \; n \in \N\}\).\\
   Si la suite \((u_n)\) est \textbf{croissante} et \textbf{non majorée} alors elle tends vers \(+ \infty\).\\
   Si la suite \((u_n)\) est \textbf{décroissante} et \textbf{non minorée} alors elle tends vers \(- \infty\).
}
Si \((u_n)\) et \((w_n)\) tendent vers une limite finie \(l \in K\), alors on a aussi le \textbf{théorème d'encadrement}:
\customBox{width=7cm}{
    \(\Bigr[\forall n \in \N \; ; \; u_n \leq v_n \leq w_n \Bigr] \implies v_n \rightarrow l\)
}

On appelle \textbf{suites adjacentes} deux suites \((u_n)\) et \((v_n)\) telles que l'une est \textbf{croissante} et l'autre \textbf{décroissante} et telles que \((u_n - v_n) \rightarrow 0\), alors on peut montrer que les deux suites convergent vers une même limite \(l \in K\) grâce aux théorèmes ci-dessus.

\subsection*{\subsecstyle{Suites complexes {:}}}
Dans le cas d'une suite à valeurs dans \(\C\), à défaut de pouvoir parler de suites majorées, on peut néanmoins parler de suites \textbf{bornées}.\<

Les suites à valeur dans \(\C\) vérifient la propriété fondamentale:
\customBox{width=5cm}{
    \(u_n = \Re(u_n) + i\Im(u_n)\)
}

Une suite à valeur dans \(\C\) converge si et seulement si \(\Re(u_n)\) et \(\Im(u_n)\) convergent.

\subsection*{\subsecstyle{Limite des bornes {:}}}
Supposons \((u_n)\) \textbf{bornée}, on construit alors les suites numériques \((a_N)\) et \((b_N)\) suivantes:
\begin{align*}
    a_N := &\sup\Bigl\{ u_n \; ; \; n > N \Bigl\} \\
    b_N := &\inf\Bigl\{ u_n \; ; \; n > N \Bigl\}
\end{align*}
\begin{center}
    \textit{
        Ici \((a_N)\) est simplement la borne supérieure de \((u_n)\) \textbf{qui a été privée des \(N\) premiers termes}.
    }
\end{center}
On peut alors remarquer que ces suites sont monotones et bornées, donc elles convergent nécessairement\footnote[1]{En effet si on prive \((u_n)\) de termes, sa borne supérieure ne peut que décroître, et \(u_n\) étant bornée, ces suite le sont aussi.} vers des limites qu'on appelle alors les \textbf{limites supérieure et inférieure} de \((u_n)\).

\begin{center}
    \textit{
        La limite supérieure d'une suite est \textbf{la plus petite borne supérieure} de cette suite quand on la prive de beaucoup de premiers termes.\+
        La limite inférieure d'une suite est \textbf{la plus grande borne inférieure} de cette suite quand on la prive de beaucoup de premiers termes.\+
    }
\end{center}

\subsection*{\subsecstyle{Comparaison Asymptotique {:}}}

On dit que la suite \((u_n)\) est \textbf{dominée} par la suite \((v_n)\) et on note alors \(u_n = O(v_n)\) si et seulement si il existe une suite \textbf{bornée} \((\theta_n)\) telle qu'à partir d'un certain rang, on ait la propriété ci dessous.\+
C'est une relation de \textbf{préordre} sur les suites, et elle est aussi \textbf{linéaire et multiplicative}.
\[
    u_n = \theta_n v_n
\]

On dit que la suite \((u_n)\) est \textbf{négligeable} devant la suite \((v_n)\) et on note \(u_n = o(v_n)\) si et seulement si il existe une suite \((\epsilon_n)\) qui tende vers 0, telle qu'à partir d'un certain rang, on ait la propriété ci dessous.\+
C'est une relation \textbf{transitive, linéaire et multiplicative}.
\[
    u_n = \epsilon_n v_n
\]

On dit que la suite \((u_n)\) est \textbf{équivalente} à la suite \((v_n)\) et on note \(u_n \sim v_n\) si et seulement si il existe une suite \((\eta_n)\) qui tende vers 1, telle qu'à partir d'un certain rang, on ait la propriété ci-dessous.\+ 
C'est une \textbf{relation d'équivalence} sur les suites, et elle est aussi \textbf{multiplicative}. 
\[
    u_n = \eta_n v_n
\]
De manière générale, on caractérise ces relations par le tableau suivant:
\begin{center}
    \renewcommand{\arraystretch}{1.5}%
    \setlength\arrayrulewidth{0.8pt}
    \begin{tabular}{| c | c | c |}
    \hline
    \(u_n = O(v_n) \Longleftrightarrow \frac{u_n}{v_n} \longrightarrow C\) &    \(u_n = o(v_n) \Longleftrightarrow \frac{u_n}{v_n} \longrightarrow 0\) &    \(u_n \sim v_n \Longleftrightarrow \frac{u_n}{v_n} \longrightarrow 1\)\\ [1.5ex]
    \hline
    \end{tabular}
 \end{center}  

\chapter*{\chapterstyle{III --- Topologie élémentaire de \(\R\)}} % 95% Fini
\addcontentsline{toc}{section}{Topologie de \(\R\)}

On admettra que \(\R\) est un \textbf{espace topologique}, on appelle \textbf{espace métrique} tout espace topologique muni d'une \textbf{distance}. En particulier, \((\R, |\cdot|)\) est alors un espace métrique.\<

Par ailleurs, on appelera \textbf{ouverts} de \(\R\) les intervalles ouverts et \textbf{fermés} les intervalles fermés.

\subsection*{\subsecstyle{Voisinages {:}}}
On appelle \textbf{voisinage} d'un point \(a \in \R\) toute partie \(V\) de \(\R\) telle que pour un certain \(\epsilon\) on ait \(\ioo{a - \epsilon}{a + \epsilon} \subset V\).\+
On appelle \textbf{voisinage} de \(+\infty\) toute partie \(V\) de \(\R\) telle que pour un certain \(A\) on ait \(\ioo{A}{+\infty} \subset V\).\+
On appelle \textbf{voisinage} de \(-\infty\) toute partie \(V\) de \(\R\) telle que pour un certain \(A\) on ait \(\ioo{-\infty}{A} \subset V\).\+
\begin{center}
   \textit{
      Un voisinage d'un point est donc une partie qui contient un ouvert qui contient ce point. Grossièrement, on peut généralement se réprésenter un voisinage comme un de ces ouverts.
   }
\end{center}
On notera \(\mathscr{V}_a\) l'ensemble des voisinages du point \(a\).\+
\underline{Exemple:} \(\ico{-2}{3}\) est un voisinage de \(3\) et de \(0\).

\subsection*{\subsecstyle{Adhérence {:}}}
On appelle \textbf{adhérence} d'une partie \(A\) de \(\R\) le plus petit fermé contenant \(A\). C'est un \textbf{opérateur de clôture} pour l'inclusion, ie c'est une application \textbf{idempotente, croissante, et extensive}.\<

Formellement, si on pose \((\mathcal{F}_i)_{i \in \N}\) la famille de tout les fermés contenant \(A\), on a:
\[
    \overline{A} := \bigcap_{i \in \N} \mathcal{F}_i
\]
Dans un \textbf{espace métrique}, on a aussi la caractérisation séquentielle:
\customBox{width=15cm}{
   L'adhérence est l'ensemble de toutes les limites finies de suites à valeurs dans \(A\).
}

\subsection*{\subsecstyle{Intérieur {:}}}
On appelle \textbf{intérieur} d'une partie \(A\) de \(\R\) le plus grand ouvert contenu dans \(A\). C'est une application \textbf{idempotente et croissante}.\<

Formellement, si on pose \((\mathcal{O}_i)_{i \in \N}\) la famille de tout les ouverts contenus dans \(A\), on a:
\[
    \accentset{\circ}{A} := \bigcup_{i \in \N} \mathcal{O}_i
\]

\subsection*{\subsecstyle{Frontière {:}}}
On appelle \textbf{frontière} d'une partie \(A\) de \(\R\) et on note \(\partial A\) la partie qui vérifie une des conditions équivalentes ci-dessous:
\[
        \partial A = \overline{A} \backslash \accentset{\circ}{A}
\]
\begin{center}
    \textit{
        La frontière est donc l'ensemble des points situés au ``bord'' de cet ensemble.
    }
\end{center}

\subsection*{\subsecstyle{Densité {:}}}
Soit \(A \subseteq \R\), alors on dit que \(A\) est \textbf{dense dans} \(\R\) si et seulement si elle vérifie une des conditions équivalentes ci-dessous:
\customBox{width=10cm}{
   Tout \textbf{ouvert} de \(\R\) contient des éléments de \(A\).\\
   L'adhérence de \(A\) est égale à \(\R\).
}


En particulier, on a par exemple les exemples canoniques suivants:
\customBox{width=7cm}{
   \begin{align*}
      &\text{\(\Q\) est dense dans \(\R\)}.\\
      &\text{\(\R \backslash \Q\) est dense dans \(\R\)}.  
  \end{align*}
}

\subsection*{\subsecstyle{Extractions {:}}}
On appelle extraction toute fonction \(\varphi : \N \rightarrow \N\) \textbf{strictement croissante}.\+
On dit que \(v_n\) est une \textbf{suite extraite} de \(u_n\) si il existe une extraction \(\varphi\) telle que:
\[
    v_n = u_{\varphi(n)}
\]
On note aussi qu'une composée d'extraction est encore une extraction.

\subsection*{\subsecstyle{Valeurs d'adhérence {:}}}
On appelle \textbf{valeur d'adhérence} tout limite d'une suite-extraite de \((u_n)\) et on peut caractériser ces valeurs d'adhérence comme suit:
\[
    \forall \epsilon \in \R_+ \; , \; \forall N \in \N \; , \; \exists n > N \; ; \; |u_n - l| < \epsilon 
\]
On peut alors montrer la propriété fondamentale:
\[
    \Bigr[u_n \rightarrow l \in \overline{K} \Bigr] \implies \Bigr[ \forall \varphi \; ; \; u_{\varphi(n)} \rightarrow l \Bigr]
\]
On montre aussi le \textbf{théorème de Bolzano-Weirstrass}:
\customBox{width=11cm}{
   Tout suite bornée admet au moins une valeur d'adhérence.
}

\subsection*{\subsecstyle{Complétude {:}}}
On appelle \textbf{suite de Cauchy} toute suite qui vérifie:
\customBox{width=9cm}{
   \(\forall \epsilon \in \R_+^* \; , \; \exists N \in \N \; , \; \forall p, q > N \; ; \; |u_p - u_q| < \epsilon \)
}

\begin{center}
    \textit{
        Une suite de Cauchy est donc une suite telle que les termes de la suite sont très proches \textbf{les uns des autres} à partir d'un certain rang\footnote[1]{Un corollaire immédiat est que \textbf{toute suite convergente est de Cauchy}.}.
    }
\end{center}

Si une suite est de Cauchy, alors elle vérifie les propriétés suivantes:
\begin{align*}
    &\bullet \text{Elle est bornée.} \\
    &\bullet \text{Si elle admet une valeur d'adhérence, alors elle converge.}
\end{align*}
\customBox{width=10cm}{
   On appelle \textbf{espace complet} tout espace métrique dans lequel \textbf{toute suite de Cauchy} converge.\<
    
   En particulier \(\R\) et \(\C\) sont des espaces complets.
}

\chapter*{\chapterstyle{III --- Limites}} % 99% Fini
\addcontentsline{toc}{section}{Limites}
On rapelle qu'un \textbf{voisinage} du point \(a\) est une partie qui contient \textbf{un ouvert} qui contient ce point et on note \(\mathscr{V}_a\) l'ensemble des voisinages du point \(a\).

\subsection*{\subsecstyle{Définition{:}}}
\addcontentsline{toc}{subsection}{Définition}

Soit \(f : D \mapsto \R\) une fonction, \(a \in \overline{D}\) et \(l \in \overline{\R}\), on dit que \(f\) admet \(l\) pour limite en \(a\) si et seulement si:
\customBox{width=6cm}{
   \(\forall \; V_l \in \mathscr{V}_l \; , \; \exists  V_a \in \mathscr{V}_a \; ; \; f(V_a) \subset V_l \)
}

Graphiquement, cela signifie simplement que peu importe l'ouvert \(O\) qu'on considère qui contient la limite, il sera possible de trouver un ouvert qui contient \(a\) tel que \textbf{tout ses points} soient envoyés dans \(O\).

\begin{center}
   \begin{tikzpicture}[domain=0:sqrt(5), yscale=0.8, xscale=1.5]
      \draw[color=DarkBlue1, thick] plot (\x,{\x^2}) node[right] {$f(x) = x^2$};
      \draw[-latex] (0,0) -- (3,0) node [right] {$x$};
      \draw[-latex] (0, 0) -- (0,5) node [above] {$y$};
      \draw[] (1.75, -0.1) -- (1.75,0.1) node[] at (1.75, -0.3) {$a$};
      \draw[] (-0.1, 1.75*1.75) -- (0.1,1.75*1.75) node[] at (-0.3, 1.75*1.75) {$l$};
      \draw[dashed] (1.75, 0.2) -- (1.75, 1.75*1.75) -- (0.14, 1.75*1.75);
   
      \fill[BrightBlue1!20,nearly transparent] (0,2) -- (3,2) -- (3,4) -- (0,4) -- cycle;
      \draw[color=BrightBlue1, thick, dashed] (0,2) -- (3,2);
      \draw[color=BrightBlue1, thick, dashed] (0,4) -- (3,4);
   
      \draw[color=BrightBlue1, line width=0.5mm] (0,2) -- (0,4);
      \draw[color=BrightBlue1, line width=0.5mm] (-0.15,2) -- (0.15,2);
      \draw[color=BrightBlue1, line width=0.5mm] (-0.15,4) -- (0.15,4);
      \node[color=BrightBlue1] at (-0.5, 4) {\(\forall V_l\)};

      \draw[color=BrightRed1, thick, dashed] (1.5,0) -- (1.5,1.5*1.5);
      \draw[color=BrightRed1, thick, dashed] (1.95,0) -- (1.95,1.95*1.95);

      \draw[color=BrightRed1, line width=0.5mm] (1.5, 0) -- (1.95, 0);
      \draw[color=BrightRed1, line width=0.5mm] (1.5, -0.15) -- (1.5, 0.15);
      \draw[color=BrightRed1, line width=0.5mm] (1.95, -0.15) -- (1.95, 0.15);
      \node[color=BrightRed1] at (1.1, -0.3) {\(\exists V_a\)};
   
      \draw[color=BrightRed1, domain=1.5:1.95, thick] plot (\x,{\x^2});
   \end{tikzpicture}
\end{center}
La nature de l'ensemble des nombres réels nous permet de caractériser cette définition et d'en déduire que \(f\) tends vers \(l\) en \(a\) si et seulement si:
\customBox{width=10.2cm}{
   \(\forall \epsilon \in \R^+ \; , \; \exists \delta \in \R^+ \; , \; \forall x \in D \; ; \; |x - a| < \delta \implies |f(x) - l| < \epsilon\)
}

Si la limite est en \(+ \infty\) (analogue pour le cas \(-\infty\)) on obtient la caractérisation suivante:
\customBox{width=10cm}{
   \(\forall \epsilon \in \R^+ \; , \; \exists M \in \R \; , \; \forall x \in D \; ; \; x > M \implies |f(x) - l| < \epsilon\)
}

Enfin si la limite est \(+ \infty\) (analogue pour le cas \(-\infty\)) on obtient:
\customBox{width=10cm}{
   \(\forall M \in \R \; , \; \exists \delta \in \R^+ \; , \; \forall x \in D \; ; \; |x - a| < \delta \implies f(x) > M\)
}

\subsection*{\subsecstyle{Caractérisation séquentielle{:}}}
\addcontentsline{toc}{subsection}{Caractérisation séquentielle}

Soit \(f\) une fonction, alors elle tends vers \(l\) en \(a\) si et seulement si pour toute suite \((u_n)_{n \in \N}\) qui tends vers \(a\) on a:
\customBox{width=3cm}{
   \(f(u_n) \underset{+\infty}{\longrightarrow} l\)
}
\begin{center}
   \textit{
      En particulier, si on arrive à trouver deux suites qui tendent vers \(a\) telles que \(f(u_n)\) et \(f(v_n)\) tendent vers deux valeurs différentes, alors la limite de \(f\) ne peut pas exister.
   }
\end{center}
\pagebreak

\subsection*{\subsecstyle{Propriétés{:}}}
\addcontentsline{toc}{subsection}{Propriétés}

Gràce aux définitions ci-dessus, on peut en déduire plusieurs propriétés, par exemple on voit directement que si \(f\) admet une limite \textbf{finie} en \(a\), alors elle est \textbf{bornée au voisinage de \(a\)}\footnote[1]{Il suffit d'utiliser la définition pour un \(\epsilon\) fixé}. \<

Les limites de fonctions héritent des mêmes propriétés opératoires que les limites de suites, en particulier, dans le cas de fonctions à limites finies, le passage à la limite est \textbf{linéaire et multiplicatif}.\+
Dans le cas des limites inifines et des quotients, des formes indéterminées peuvent apparaître et nécessitent les précautions usuelles pour conclure.\<

Cette définion, aussi appellée \textbf{limite pointée}, nous permet alors d'avoir le \textbf{théorème de composition de limites}.\+
En effet, pour \(a, b, l \in \overline{\R}\) et deux fonctions telles que \(f \underset{b}{\rightarrow}l\) et \(g \underset{a}{\rightarrow}b \) on a:
\customBox{width=4cm}{
   \(\lim_{a} (f \circ g)(x) = l\)
}

Enfin, tout les théorèmes d'encadrements des suites s'étendent dans le cadre des fonctions.\<

On peut considérer les restrictions de \(f\) à droite ou à gauche de \(a\) et ainsi définir \textbf{les limites latérales} de \(f\) en \(a\) qui sont analogues à la définition principale.

\subsection*{\subsecstyle{Comparaison asymptotique II{:}}}
\addcontentsline{toc}{subsection}{Comparaison asymptotique II}

Les trois différentes relations de comparaison se comportent de la même manière avec les fonctions qu'avec les suites. En particulier, on a toujours:
\begin{center}
   \renewcommand{\arraystretch}{1.5}%
   \setlength\arrayrulewidth{0.8pt}

   \begin{tabular}{| c | c | c |}
   \hline
   \(f \underset{a}{=} O(g) \Longleftrightarrow \frac{f}{g} \underset{a}{\longrightarrow} C\) &    \(f \underset{a}{=} o(g) \Longleftrightarrow \frac{f}{g} \underset{a}{\longrightarrow} 0\) &    \(f \underset{a}{\sim} g \Longleftrightarrow \frac{f}{g} \underset{a}{\longrightarrow} 1\)\\ [1.5ex]
   \hline
   \end{tabular}
\end{center}  
On peut notamment caractériser le fait qu'une limite existe par:
\customBox{width=4.5cm}{
   \(f \longrightarrow l \Longleftrightarrow f = l + o(1)\)
}
Aussi, une caractérisation trés utile des équivalents nous donne\footnote[2]{La démonstration est triviale et utilise simplement la caractérisation en terme de quotient}:
\customBox{width=5cm}{
   \(f \sim g \Longleftrightarrow f = g + o(g)\)
}

Enfin, il est utile de connaître ces équivalents usuels, on considère le cas général d'une fonction \(u\) qui tends vers 0 en \(a\), alors on a les équivalents suivants (en \(a\)):

\begin{center}
   \renewcommand{\arraystretch}{1.5}%
   \setlength\arrayrulewidth{0.8pt}

   \begin{tabular}{| c | c | c |}
   \hline
   \(ln(1 + u) \sim u\) & \(e^u - 1 \sim u\) & \((1 + u)^\alpha - 1 \sim \alpha u\)\\ [0.5ex]
   \hline
   \(\sin(u) \sim u\) & \(\cos(u) \sim 1\) & \(\tan(u) \sim u\)\\ [0.5ex]
   \hline
   \end{tabular}
\end{center} 
\begin{center}
   \textit{
      On remarquera dans le chapitre sur les développements limités que la grande majorité de ces équivalents ne sont que le premier terme non-nul du développement limité d'une fonction.
   }
\end{center}

\chapter*{\chapterstyle{III --- Continuité}} % 99% Fini
\addcontentsline{toc}{section}{Continuité}

On rapelle qu'une partie de \(\R\) est \textbf{connexe} si elle ne peut pas être écrite comme réunion d'ouverts disjoints, informellement, c'est un ensemble "d'un seul tenant".\+
On rapelle qu'une partie de \(\R\) est \textbf{compacte} si c'est une partie fermée et bornée, typiquement les segements ou les unions de segments.

\subsection*{\subsecstyle{Définition{:}}}
\addcontentsline{toc}{subsection}{Définition}

On dit que la fonction \(f\) est \textbf{continue} en \(a\) si et seulement si:
\customBox{width=3.5cm}{
   \(f(x) \underset{a}{\longrightarrow} f(a)\)
}
\begin{center}
   \textit{
      Une fonction est donc continue en un point \(a\) si sa limite en \(a\) est égale à sa valeur évaluée en \(a\).
   }
\end{center}
Le lien clair entre le concept de limite et celui de continuité nous permet de définir \textbf{la continuité latérale} d'une fonction, qui est simplement la continuité à droite ou à gauche qui découle directement de notre définition des limites latérales d'une fonction.\<

On peut étendre cette définition à la continuité sur \textbf{un intervalle} \(I\), en la définissant comme la continuité en \textbf{tout points de \(I\)}.\+

On note \(\mathscr{C}(D, \K)\) l'ensemble des fonctions continues sur \(D\) à valeurs dans \(\K\).

\subsection*{\subsecstyle{Propriétés{:}}}
\addcontentsline{toc}{subsection}{Propriétés}

Une des propriétés fondamentales des fonctions continues et de l'ensemble \(\mathscr{C}(D, \K)\) est qu'il est stable par somme, multiplication externe et mulplication interne.
\begin{center}
   \textit{
      En d'autres termes cet ensemble est non seulement un \textbf{sous-espace vectoriel} de l'espace des fonctions, mais même une \textbf{sous-algèbre} de cet espace.
   }
\end{center}
De plus, l'inverse d'une fonction continue qui ne s'annule pas est aussi une fonction continue, et la composition de fonctions continues composables est aussi une fonction continue\footnote[1]{
   Attention pour la composition il faut bien considèrer les domaines de défintion manipulés.}.

\subsection*{\subsecstyle{Caractérisation séquentielle{:}}}
\addcontentsline{toc}{subsection}{Caractérisation séquentielle}

Comme dans le cas des limites on peut caractériser la continuité par les suites, en effet \(f\) est continue en \(a\) si et seulement si pour toute suite \((u_n)_{n\in\N}\) qui admet pour limite \(a\) on a:
\customBox{width=3.5cm}{
   \(f(u_n) \underset{+\infty}{\longrightarrow} f(a)\)
}

\subsection*{\subsecstyle{Prolongement par continuité {:}}}
\addcontentsline{toc}{subsection}{Prolongement par continuité}

Considérons une fonction \(f\) continue sur un intervalle \(I \backslash \{a\}\), telle que sa limite en \(a\) existe, alors on peut considèrer la fonction \(\widetilde{f}\) définie par:
\[
   \begin{aligned}
      \widetilde{f}: I &\longrightarrow \R \; ; \; x &\longmapsto \begin{cases}
         \widetilde{f}(a) = \lim_{a} f(x) \text{ si } x = a\\
         \widetilde{f}(x) = f(x) \text{ sinon}
      \end{cases}
   \end{aligned}
\]
Alors par construction \(\widetilde{f}\) est bien continue sur tout \(I\) et on l'appelle \textbf{le prolongement par continuité} de \(f\) en \(a\).

\subsection*{\subsecstyle{Théorèmes généraux{:}}}
\addcontentsline{toc}{subsection}{Théorèmes généraux}

De la continuité découlent des grands théorèmes généraux dans l'études des fonctions, en premier lieu on déduit le \textbf{théorème des valeurs intérmédiaires} et on peut alors montrer que:
\customBox{width=12cm}{
   L'image d'un \textbf{connexe} par une fonction continue est un \textbf{connexe}.
}
Ou dans sa forme analytique pour \(f\) une fonction continue sur \(\icc{a}{b}\):
\customBox{width=12cm}{
   Pour tout réel entre \(f(a)\) et \(f(b)\), il en existe un antécédent entre \(a\) et \(b\).
}

Remarquer que l'antécédent en question \textbf{n'est pas nécessairement unique.}. Graphiquement, on a:
\begin{center}
   \begin{tikzpicture}[domain=0:sqrt(5), yscale=0.8, xscale=1.5]
      \draw[color=DarkBlue1, thick] plot (\x,{\x^2}) node[right] {$f(x) = x^2$};
      \draw[-latex] (0,0) -- (3,0) node [right] {$x$};
      \draw[-latex] (0, 0) -- (0,5) node [above] {$y$};
      \draw[color=BrightBlue1] (1.75, -0.1) -- (1.75,0.1) node[] at (1.75, -0.55) {$c$};
      \draw node[] at (1.414, -0.55) {$a$};
      \draw node[] at (2, -0.5) {$b$};

      \draw[color=BrightBlue1] (-0.1, 1.75*1.75) -- (0.1,1.75*1.75) node[] at (-0.5, 1.75*1.75) {$f(c)$};
      \draw[dashed] (1.75, 0.2) -- (1.75, 1.75*1.75) -- (0.14, 1.75*1.75);
   
      \fill[BrightBlue1!20,nearly transparent] (0,2) -- (3,2) -- (3,4) -- (0,4) -- cycle;
      \draw[color=BrightBlue1, thick, dashed] (0,2) -- (3,2);
      \draw[color=BrightBlue1, thick, dashed] (0,4) -- (3,4);
   
      \draw node[] at (-0.5, 4) {$f(b)$};
      \draw node[] at (-0.5, 2) {$f(b)$};

      \draw[color=BrightBlue1, line width=0.5mm] (0,2) -- (0,4);
      \draw[color=BrightBlue1, line width=0.5mm] (-0.10,2) -- (0.10,2);
      \draw[color=BrightBlue1, line width=0.5mm] (-0.10,4) -- (0.10,4);

      \draw[color=BrightBlue1, thick, dashed] (1.414,0) -- (1.414,2);
      \draw[color=BrightBlue1, thick, dashed] (2,0) -- (2,4);

      \draw[color=BrightBlue1, line width=0.5mm] (1.414, 0) -- (2, 0);
      \draw[color=BrightBlue1, line width=0.5mm] (1.414, -0.15) -- (1.414, 0.15);
      \draw[color=BrightBlue1, line width=0.5mm] (2, -0.15) -- (2, 0.15);
   \end{tikzpicture}
\end{center}

On en déduit \textbf{le théorème de la bijection}\footnote[1]{Voir la partie ci-aprés pour une définition du concept d'homéomorphisme.}, si \(f\) est continue sur \(\ioo{a}{b}\) et \textbf{strictement monotone} alors:
\customBox{width=10cm}{
   Elle réalise \textbf{un homéomorphisme sur son image}.
}

On montre aussi le \textbf{thèorème des bornes atteintes}, si \(f\) est continue, alors on a:
\customBox{width=12cm}{
   L'image d'un \textbf{compact} par une fonction continue est un \textbf{compact}.
}

En particulier l'image d'un segment par une fonction continue est aussi un segment.

\subsection*{\subsecstyle{Continuité uniforme {:}}}
\addcontentsline{toc}{subsection}{Continuité uniforme}

Il existe une propriété de régularité \textbf{plus forte que la continuité} appelée \textbf{continuité uniforme}\footnote[2]{Il est important de noter que la continuité uniforme n'est plus une propriété locale, mais globale, sur un intervalle.} définie comme:
\customBox{width=12cm}{
   \(\forall \epsilon > 0 \; , \; \exists \delta_{\epsilon} > 0  \; , \; \forall (x_1, x_2) \in I^2 \; ; \; |x_1 - x_2| < \delta_{\epsilon} \implies |f(x_1) - f(x_2)| < \epsilon \)
}
La différence fondamentale entre ces deux propriétés étant que \(\delta\) \textbf{ne dépend plus que de \(\epsilon\)}. En particulier on peut l'interpréter:
\begin{center}
   \textit{Pour une distance \(\epsilon\) donnée, si deux réels sont assez proches (au sens \(\delta\)-proches), alors leurs images seront assez proches (au sens \(\epsilon\)-proches).}
\end{center}

Un des théorèmes principaux en relation avec la continuité uniforme est \textbf{le théorème de Heine} qui montre:
\customBox{width=13cm}{
   Tout fonction continue sur un \textbf{compact} y est \textbf{uniformément continue}.
}
En particulier, tout fonction continue sur un segment y est uniformément continue.
\pagebreak

\subsection*{\subsecstyle{Lipschitzianité{:}}}
\addcontentsline{toc}{subsection}{Lipschitzianité}

Il existe une condition de régularité \textbf{encore plus forte} que la continuité uniforme qui est le caractère \(K\)-lipschitzien. On dit qu'une fonction est \(K\)-lipschitzienne pour \(K \in \R\) si et seulement si pour tout \(x, y\) distincts dans le domaine de définition, on a:
\customBox{width=5cm}{
   \(|f(x) - f(y)| \leq K|x - y|\)
}
\begin{center}
   \textit{
      Une telle fonction a donc la particularité d'avoir \textbf{ses pentes bornées} par un réel fixé\footnote[1]{Il suffit de diviser par \(|x - y|\) non nul pour le voir.}.
   }
\end{center}
La \(K\)-lipschitziannité étant une propriété plus forte que la continuité uniforme, plus précisément, on a la suite d'implications:
\customBox{width=12cm}{
   \(K\)-lipschitzienne \(\implies\) Uniformément continue\(\implies\) Continue
}

Par ailleurs si \(K \in \ico{0}{1}\), alors on dit que \(f\) est \textbf{contractante}, ce type de fonctions a souvent un trés bon comportement de convergence quand on le retrouve dans des définitions de suites par récurrence.

\subsection*{\subsecstyle{Homéomorphismes {:}}}
\addcontentsline{toc}{subsection}{Homéomorphismes}

On peut généraliser le concept de fonctions continues dans un cadre plus général d'applications entre structures, en particulier entre \textbf{espaces topologiques}. On utilise alors une définition différente de la continuité, une définition \textbf{globale}:
\begin{center}
   \textit{L'image réciproque d'un ouvert est un ouvert\footnote[2]{Les ouverts sont les éléments même d'une topologie.}.}
\end{center}

Soit \(\varphi\) une application entre deux tels espaces telle que \(\varphi\) soit bijective, continue et \textbf{dont la réciproque est continue}, alors \(\varphi\) est un \textbf{homéomorphisme} et on dit que les deux espaces sont \textbf{homéomorphes}.
\begin{center}
   \textit{
      C'est une application fondamentale du concept de continuité car les homémorphismes permettent exactement de caractériser des espaces identiques à déformation continue près\footnote[3]{Formellement, les homéomorphismes sont exactement les isomorphismes d'espaces topologiques.}.
   }
\end{center}
\underline{Exemple:} Du point de vue de la topologie, un donut est une tasse c'est à dire que l'on peut déformer un donut continument pour obtenir une tasse.

\chapter*{\chapterstyle{III --- Dérivabilité}} % 99% Fini
\addcontentsline{toc}{section}{Dérivabilité}

\subsection*{\subsecstyle{Définition {:}}}
\addcontentsline{toc}{subsection}{Définition}

On dit qu'une fonction \(f\) est dérivable en un point \(a\) si et seulement si la limite ci-dessous existe et est \textbf{finie} et on définit alors le nombre dérivé \(f'(a)\):
\customBox{width=4.5cm}{
   \[f'(a) := \underset{x \rightarrow a}{\lim} \frac{f(x) - f(a)}{x - a}\]
}
On peut alors définir la dérivée de la fonction \(f\) qu'on note \(f'\)  qui est la fonction qui à chaque point où \(f\) est dérivable lui associe son nombre dérivé, ainsi que \textbf{la tangente à la courbe} de \(f\) en \(a\) comme étant la droite:
\customBox{width=5cm}{
   \(y := f(a) + f'(a)(x - a)\)
}

A nouveau, le lien clair entre le concept de limite et celui de dérivabilité nous permet de définir \textbf{la dérivabilité latérale} d'une fonction, qui est simplement la dérivabilité à droite ou à gauche qui découle directement de notre définition des limites latérales d'une fonction.\<

On peut étendre cette définition à la dérivabilité sur \textbf{un intervalle} \(I\), en la définissant comme la dérivabilité en \textbf{tout points de \(I\)}.\+
On note \(\mathscr{D}(D, \K)\) l'ensemble des fonctions dérivables sur \(D\) à valeurs dans \(\K\), par ailleurs il est important de noter que \textbf{la dérivabilité implique la continuité}, ie formellement on a \(\mathscr{D}(D, \K) \subset \mathscr{C}(D, \K)\).

\subsection*{\subsecstyle{Propriétés {:}}}
\addcontentsline{toc}{subsection}{Propriétés}

Une des propriétés fondamentales des fonctions dérivables et de l'ensemble \(\mathscr{D}(D, \K)\) est qu'il est stable par somme, multiplication externe et mulplication interne.
\begin{center}
   \textit{
      En d'autres termes cet ensemble est non seulement un \textbf{sous-espace vectoriel} de l'espace des fonctions (continues), mais même une \textbf{sous-algèbre} de cet espace.
   }
\end{center}
De plus, l'inverse d'une fonction dérivable qui ne s'annule pas est aussi une fonction dérivable, et la composition de fonctions dérivables composables est aussi une fonction dérivable.

\subsection*{\subsecstyle{Etudes des variations {:}}}
\addcontentsline{toc}{subsection}{Etudes des variations}

Si on considère une fonction dérivable, alors on appelle \textbf{point critique} tout point \(a\) où la fonction est dérivable et où \(f'(a) = 0\). Intuitivement, c'est un point du domaine de définition tel que la fonction ne varie pas en ce point.\+
En particulier, une fonction est dérivable en un point \textbf{intérieur}\footnote[1]{Si le point n'est pas intérieur et est un extremum, la dérivée n'est pas nécessairement nulle. Par exemple la restriction (à l'arrivée) de ln(x) à \(\R^-\) est dérivable à gauche en 1, y admet un maximum, mais la dérivée n'est pas nulle.} et y admet un \textbf{extremum local}, alors c'est un point critique.\<

Enfin le signe de la dérivée nous permet de comprendre les variations de la fonction sur un \textbf{intervalle}, en particulier on a la propriété suivante:
\begin{center}
   \textit{Si \(f'\) conserve le même signe sur un intervalle, alors \(f\) est monotone sur cet intervalle\footnote[2]{En particulier elle est croissante dans le cas positif et décroissante sinon}.}
\end{center}
En particulier, si \(f'\) est nulle sur un intervalle, alors \(f\) est constante sur cet intervalle.
\pagebreak 

\subsection*{\subsecstyle{Classes de régularité {:}}}
\addcontentsline{toc}{subsection}{Classes de régularité}

Les classes de régularité des fonctions numériques constituent une classification des fonctions basées sur l'existence et la continuité des dérivées itérées.
Par exemple on note \(\mathscr{C}^0(D, \K)\) l'ensemble des fonctions continues, \(\mathscr{C}^1(D, \K)\) l'ensemble des fonctions dérivables à dérivée continue.\<

En généralisant ces notations, on note \(\mathscr{C}^k(D, \K)\) l'ensemble des fonctions \(k\) fois dérivables \textbf{et à dérivée continue}.\<

Ces ensembles \textbf{héritent des propriétés de stabilité} (par somme, produit interne et externe, composition et inverses) des ensembles de base définis précédemment. En particulier ce sont tous des sous-algèbres de l'espace des fonctions.

\subsection*{\subsecstyle{Dérivation {:}}}
\addcontentsline{toc}{subsection}{Dérivation}

Soit \(f\) une fonction dérivable, alors on peut appliquer à \(f\) l'opérateur de dérivation, en d'autres termes calculer sa dérivée \(f'\). Dans cette partie nous allons détailler les propriétés élémentaires de cet opérateur:
\begin{center}
   \renewcommand{\arraystretch}{1.5}%
   \setlength\arrayrulewidth{0.8pt}

   \begin{tabular}{| c | c | c |}
   \hline
   \((f + g)' = f' + g'\) & \((fg)' = f'g + fg'\)  & \((f \circ g)' = (f' \circ g) \cdot g'\) \\ [1ex]
   \hline
   \end{tabular}
\end{center}  
Soit \(f\) de classe \(\mathscr{C}^n\), alors dans le cas de dérivées successives, la propriété pour la somme ci dessus se généralise aisément en:
\[
   (f + g)^{(n)} = f^{(n)} + g^{(n)}
\]
Et la formule du produit, appellée \textbf{formule de Liebniz} se généralise aussi en:
\customBox{width=5.5cm}{
   \[(fg)^{(n)} = \sum_{k=0}^{n}\binom{k}{n}f^{(k)}g^{(n - k)}\]
}
Par ailleurs, si on pose \(f\) une fonction dérivable sur \(I\) et que \(f'\) ne s'annule pas sur \(I\), alors la réciproque de \(f\) est dérivable\footnote[1]{Les deux hypothèses impliquent que \(f\) est bijective sur son image, puis on revient à la définition pour montrer que \(f^{-1}\) est dérivable en \(f(a)\), on peut effectuer un changement de variable dans le calcul de limite pour trouver la première identité. La seconde nécessite un second changement de variable évident.} pour tout \(y := f(a)\) sur cet intervalle et on a:
\customBox{width=6cm}{
   \[(f^{-1})'(y) = \frac{1}{f'(a)} = \frac{1}{f'(f^{-1}(y))}\]
}
On peut généraliser ce résultat\footnote[2]{Même début de preuve que ci-dessus, avec une récurrence pour conclure.} en montrant que si \(f\) est une fonction de classe \(\mathscr{C}^k\) et si \(f'\) \textbf{ne s'annule pas}, alors \(f^{-1}\) est aussi de classe \(\mathscr{C}^k\).\<

Cette propriété permet de prouver la dérivabilité d'une réciproque, ainsi que de trouver sa dérivée.\+
\underline{Exemple:} \(f: x \mapsto x^2\) est dérivable et sa dérivée ne s'annule pas sur cet \(\icc{1}{2}\), donc elle est bijective sur \(\icc{1}{4}\) et sa réciproque est dérivable de dérivée \(\frac{1}{2\sqrt{x}}\).

\subsection*{\subsecstyle{Accroissements finis {:}}}
\addcontentsline{toc}{subsection}{Accroissements finis}

Soit \(f\) une fonction continue sur un intervalle \(\icc{a}{b}\) et dérivable sur \(\ioo{a}{b}\), le \textbf{théorème des accroissements finis} assure l'existence d'un \(c \in \ioo{a}{b}\) tel que:
\customBox{width=4.5cm}{
   \[\frac{f(b) - f(a)}{b - a} = f'(c)\]
}
Intuitivement, on comprends alors que la quantité à droite est le taux d'acroissement global entre \(a\) et \(b\), et donc le théorème nous assure l'existence d'un point tel que \textbf{le taux d'acroissement en ce point soit égal au taux d'acroissement global}.
\pagebreak

Graphiquement, on voit sur le dessin ci-dessous que les deux taux d'accroissement sont égaux:
\begin{center}
   \begin{tikzpicture}[domain=0:4]
      \draw[color=DarkBlue1, thick] plot (\x,{3*\x - 0.75*\x^2});

      \draw[-latex] (0,0) -- (5,0) node [right] {$x$};
      \draw[-latex] (0, 0) -- (0,4) node [above] {$y$};


      \draw[] (0.5, -0.1) -- (0.5,0.1) node[] at (0.5, -0.3) {$a$};
      \draw[color=BrightRed1, thick, dashed] (0.5,0) -- (0.5, 1.3125);
      \draw[] (2.5, -0.1) -- (2.5,0.1) node[] at (2.5, -0.3) {$b$};
      \draw[color=BrightRed1, thick, dashed] (2.5,0) -- (2.5, 2.8125);

      \draw[color=BrightRed1, thick] (0.5, 1.3125) -- (2.5, 2.8125) node at (3.8, 2.9) {$\tau_1 = \frac{f(b) - f(a)}{b - a}$};

      \draw[color=BrightBlue1] (1.5, -0.1) -- (1.5,0.1) node[] at (1.5, -0.3) {$\exists c$};
      \draw[color=BrightBlue1, thick, dashed] (1.5, 0) -- (1.5, 2.8125);
      \draw[color=BrightBlue1, thick, domain=0.5:2.5] plot (\x,{1.6875 + 0.75*\x}) node[] at (1, 3.3) {$\tau_2 = f'(c)$};

   \end{tikzpicture}
\end{center}
Un cas particulier remarquable de ce théorème est le \textbf{théorème de Rolle} dans le cas où \(f(a) = f(b)\), donc dans le cas où le taux d'acroissement global est nul, alors on a l'existence d'un \(c \in \ioo{a}{b}\) tel que \(f'(c) = 0\).\<

Si \(f'\) est bornée\footnote[1]{En effet, la lipschitzianité est exactement le concept d'une fonction ayant ses pentes bornées.}, on montre aussi \textbf{l'inégalité des accroissements finis}:
\customBox{width=8cm}{
   \(f\) est \(K\)-lipschitzienne avec \(K = \sup|f'|\)
}
\underline{Exemple:} \(|\sin(x)'| = |\cos(x)| \leq 1\) donc sinus est une fonction \(1\)-lipschitzienne et donc on montre par exemple que \(|\sin(x) - \sin(0)| \leq |x - 0| \Longleftrightarrow |\sin(x)| \leq |x|\)

\subsection*{\subsecstyle{Théorème de la limite de la dérivée{:}}}
\addcontentsline{toc}{subsection}{Théorème de la limite de la dérivée}

Soit \(f\) une fonction continue sur \(I\) et dérivable sur \(I \; \backslash \; \{a\}\), alors on peut montrer que si \(\underset{x \rightarrow a}{\lim} f'(x) =  l\), alors \(f\) \textbf{est dérivable} en \(a\) de dérivée \(l\).\<

En particulier pour une fonction \(f \in \mathscr{C}^1(I \; \backslash \; \{a\}, \R)\), si \({\lim}f\) et \({\lim}f'\) existent, alors le prolongement à \(I\) tout entier est de classe \(\mathscr{C}^1\).\+
En général pour une fonction \(f \in \mathscr{C}^k(I \; \backslash \; \{a\}, \R)\), si les limites de toutes les dérivées intermédiaires existent, alors le prolongement à \(I\) tout entier est de classe \(\mathscr{C}^k\).

\subsection*{\subsecstyle{Difféomorphismes {:}}}
\addcontentsline{toc}{subsection}{Difféomorphismes}

A l'instar des homéomorphismes le concept de fonctions dérivables se généralise dans un cadre plus général d'applications \textbf{différentiables} entre structures, en particulier entre \textbf{variétés différentielles}. On utilise alors une définition différente de la dérivabilité.\<

Soit \(\varphi\) une application entre deux tels espaces telle que \(\varphi\) soit bijective, différentiable et \textbf{dont la réciproque est différentiable}, alors \(\varphi\) est un \textbf{difféomorphismes} et on dit que les deux espaces sont \textbf{difféomorphes}.
\begin{center}
   \textit{
      C'est une application fondamentale du concept de dérivabilité car les difféomorphismes permettent exactement de caractériser des espaces identiques du point de vue de l'analyse.
   }
\end{center}
Par exemple les difféomorphismes permettent de préserver le caractère \textbf{tangent} entre deux espaces.

\chapter*{\chapterstyle{III --- Convexité}} % 99% Fini
\addcontentsline{toc}{section}{Convexité}

\subsection*{\subsecstyle{Définition {:}}}
\addcontentsline{toc}{subsection}{Définition}

On dit que \(f\) est \textbf{convexe} sur \(I\) si et seulement si \textbf{son graphe est en dessous de toutes ses cordes}, ie elle vérifie la propriété:
\customBox{width=10cm}{
   \(\forall a, b \in I \; , \; \forall \lambda \in \icc{0}{1}\; ; \; f((1-\lambda)a + \lambda b) \leq (1-\lambda)f(a) + \lambda f(b)\)
}
Une fonction qui vérifie l'inégalité inverse est dite \textbf{concave} et son graphe est situé \textbf{au dessus de toutes ses cordes}.
Graphiquement, voici un exemple de fonction convexe:
\begin{center}
   \begin{tikzpicture}[domain=0:4, yscale=0.8]
      \draw[color=DarkBlue1, thick, dashed] plot (\x,{-3*\x + 0.75*\x^2+3.5}) node [right] {$f(x)$};

      \draw[-latex] (0,0) -- (5,0) node [right] {$x$};
      \draw[-latex] (0, 0) -- (0,4) node [above] {$y$};

      \draw[|-|, thick, color=BrightRed1] (1, 0) -- (3.5,0);
      \draw node at (1, -0.4) {$a$};
      \draw node at (3.5, -0.4) {$b$};
      \draw[dashed, color=BrightRed1, thick] (1, 0) -- (1, 1.25);      
      \draw[dashed, color=BrightRed1, thick] (3.5, 0) -- (3.5, 2.1875);

      \draw[color=BrightRed1, thick, domain=1:3.5] plot (\x,{-3*\x + 0.75*\x^2+3.5}) node [below right] {$f((1-\lambda)a + \lambda b)$};

      \draw[color=DarkBlue1, thick] (1, 1.25) -- (3.5, 2.1875) node[] at (5.35, 2.5) {$(1-\lambda)f(a) + \lambda f(b)$};

   \end{tikzpicture}
\end{center}
\begin{center}
   \textit{
      En d'autres termes l'image d'un barycentre (à coefficient positif) est en dessous du barycentre des images.
   }
\end{center}
Par ailleurs il existe une caractérisation trés utile de la convexité, pour montrer qu'une fonction \textbf{dérivable} est convexe, il suffit de montrer que sa \textbf{dérivée est croissante}, ou alors, si elle est deux fois dérivable, que sa \textbf{dérivée seconde est positive}.\<

\underline{Exemple:}  La fonction logarithme est concave sur \(\R\) car sa dérivée seconde est négative, donc pour tous \(a, b > 0\), on a \(\ln(\frac{a + b}{2}) \leq \frac{1}{2}\ln(a) + \frac{1}{2}\ln(b)\).

\subsection*{\subsecstyle{Propriétés {:}}}
\addcontentsline{toc}{subsection}{Propriétés}

On peut noter deux grandes propriétés des fonctions convexes, tout d'abord pour tout \(a, b, c\) dans \(I\) tels que \(a < b < c\), on a \textbf{l'inégalité des pentes} qui nous donne:
\customBox{width=7.5cm}{
   \[\frac{f(b)-f(a)}{b-a} \leq \frac{f(c)-f(a)}{c-a} \leq \frac{f(c)-f(b)}{c-b}\]
}
Intuitivement, cela signifie que \textbf{les pentes des cordes sont croissantes}.\<

Finalement on peut généraliser l'inégalité de la définition par \textbf{l'inégalité de Jensen}, qui pour toute famille \((x_i)_{i \in \N}\) de points et famille \((\lambda_i)_{i \in \N}\) de coefficients positifs de somme 1, nous donne:
\customBox{width=5cm}{
   \[f\biggl(\sum_{i \in \N} \lambda_i x_i\biggl) \leq  \sum_{i \in \N} \lambda_i f(x_i)\]
}
Qui généralise l'intuition selon laquelle \textbf{l'image d'un barycentre est en dessous du barycentre des images}.

\chapter*{\chapterstyle{III --- Développements limités}} % 99% Fini
\addcontentsline{toc}{section}{Développements limités}

Dans ce chapitre nous cherchons à approximer \textbf{localement} des fonctions par des polynômes au voisinage d'un point ou de l'infini.\<

Il est important de noter que la majorité des égalités de ce chapitre sont des égalités \textbf{au voisinage d'un point}, et non des égalités globales.
\subsection*{\subsecstyle{Définition {:}}}
\addcontentsline{toc}{subsection}{Définition}

Soit \(f\) une fonction de \(D\) dans \(\R\), et \(a\) un point adhérent à \(D\). On dit que \(f\) admet un \textbf{développement limité} d'ordre \(n\) \textbf{au voisinage de } \(a\) et on note \(DL_n(a)\) si il existe une famille de réels \((a_i)\) tels que:
\customBox{width=10cm}{
   \(f(x) = a_0 + a_1(x-a) + \ldots + a_n(x-a)^n + o((x-a)^n)\)
}
Si elle existe, alors cette famille est \textbf{unique}, par ailleurs plus l'ordre est grand, plus la précision est grande.\+

La partie polynomiale d'une développement limité est appellée \textbf{partie principale}.

\subsection*{\subsecstyle{Propriétés{:}}}
\addcontentsline{toc}{subsection}{Propriétés}

Soit \(f\) et \(g\) deux fonctions qui admettent un développement limité à l'ordre \(n\).\+
On appelle \textbf{troncature} à l'ordre \(k\) de \(f\), la partie principale de \(f\) privée de tout les coefficients de degré supérieurs à \(k\), et suivi d'un \(o(x^k)\). Il s'agit simplement ici de réduire la précision de notre approximation, en particulier, si on a un développement limité à l'ordre \(n\) alors on a un développement limité à tout ordre inférieur à \(n\) par troncature.\<

Aussi, les développements limités se comportent de manière intuitive par rapports aux opérations algébriques:
\begin{center}
   \textit{Pour calculer un développement limité d'une somme, il suffit de faire la somme des parties principales. \+
   Pour la multiplication (la composition), il suffit de faire le produit (la composée) des parties principales.\footnote[1]{\textbf{Attention:} Cette méthode produira dans la plupart des cas des termes inutiles, il faut donc tronquer et/ou anticiper le degré du produit final, pour faciliter les calculs, voir l'exemple.} }
\end{center}

\underline{Exemple:} Cherchons un \(DL_2(e^x\cos(x))\), on a \(e^x\cos(x) = [1 + x + \frac{x^2}{2}+o(x^2)] [1 - \frac{x^2}{2} + o(x^2)]  \) = \(1 + x + o(x^2)\)

\subsection*{\subsecstyle{Changement de voisinage {:}}}
\addcontentsline{toc}{subsection}{Changement de voisinage}

On peut ramener la recherche d'un déveoppement limité en un point à un développement limité en \(0\) par une translation.\footnote[1]{En effet, un DL\((a, f(x))\) est un DL\((0, f(x + a))\) à translation près (faire un dessin). Il suffit donc de chercher un DL\((0, f(x + a))\) puis de translater à nouveau sur la fonction originelle.}\<

On peut ramener la recherche d'un déveoppement asymptotique en \(\pm \infty\), à un développement limité en \(0\) par le changement de variable
\( X = \frac{1}{x}\).\<

\underline{Exemple:} La fonction \(e^\frac{1}{x}\) est définie au voisinage de \(+ \infty\), on pose \(X = \frac{1}{x}\), alors un \(DL_2(0)\) de \(f(X)\) est \(1 + X + \frac{X^2}{2}+o(X^2) = 1 + \frac{1}{x} + \frac{1}{2x^2} + o(\frac{1}{x^2})\) aprés substitution.
\pagebreak

\subsection*{\subsecstyle{Formules de Taylor {:}}}
\addcontentsline{toc}{subsection}{Formule de Taylor-Young}
Soit \(f\) une fonction de classe \(\mathscr{C}^n\) sur un intervalle \(I\), alors on peut montrer \textbf{l'existence} d'un développement limité à l'ordre \(n\) au voisinage d'un point \(a\) de \(I\) qui est donné par:
\customBox{width=7cm}{
   \[f(x) = \sum_{k=0}^{n} \frac{f^{(k)}(a)}{k!}(x - a)^k + o((x - a)^n)\]
}
La quantité négligeable appellée \textbf{reste}, qu'on note généralement \(R_n\), peut s'expliciter dans le cas où la fonction est de classe \(\mathscr{C}^{n+1}\) et pour un certain \(c \in \icc{a}{x}\), on a:
\customBox{width=11cm}{
   \[ R_n = \frac{f^{(k+1)}(c)}{k!}(x - a)^{k+1} 
   \quad\quad\quad\quad
   R_n = \int_{a}^{x} \frac{f^{(k+1)}(t)}{k!}(x - t)^{k} \d t\]
}  
Ce sont respectivement les formules de \textbf{Taylor-Lagrange}\footnote[1]{Celle ci peut se comprendre comme une généralisation du théorème des accroissements finis qui en est un cas particulier (à l'ordre 1).} et de \textbf{Taylor avec reste intégral}.\+
Expliciter le reste permet alors d'obtenir des informations sur tout l'intervalle \(\icc{a}{x}\), en effet, ce ne sont alors plus des développements asymptotiques mais bien des formules \textbf{exactes}.\<

\underline{Exemple:} Pour un certain \(c\) sur l'intervale \(\icc{0}{x}\) , on a \(\cos(x) = 1 - \frac{x^2}{2!} + \frac{x^4}{4!} + \frac{x^5}{5!}\sin(c)\), cela nous permet alors d'avoir, par exemple, la majoration \(\cos(x) \leq 1 - \frac{x^2}{2!} + \frac{x^4}{4!} + \frac{x^5}{5!}\).
\<


Aussi, \(f\) est continue en \(a\) si et seulement si elle admet un \textbf{développement limité à l'ordre 0}.\+
De même, elle est dérivable en \(a\) si et seulement si elle admet un \textbf{développement limité à l'ordre 1}.\+
Néanmoins, au dela de l'ordre \(k = 1\), admettre un développement limité à l'ordre \(k\) n'implique \textbf{pas} d'être de classe \(\mathscr{C}^k\).

\subsection*{\subsecstyle{Intégration d'un développement limité{:}}}
\addcontentsline{toc}{subsection}{Intégration d'un développement limité}

Une propriété trés particulière nous permet de trouver le développement limité d'une fonction si l'on connaît le développement limité de sa dérivée.\+ 
Soit \(f\) une fonction dérivable tel que \(f'\) admet un développement limité à l'ordre \(n\), ie:
\[
   f'(x) \underset{0}{=} \sum_{k=0}^{n} a_kx^k + o(x^k)   
\]
Alors \(f\) admet un développement limité à l'ordre \(n + 1\) obtenu \textbf{en intégrant termes à termes}, ie on a:
\[
   f(x) \underset{0}{=} f(a) + \sum_{k=0}^{n} \frac{a_k x^{k+1}}{k+1} + o(x^{k+1})   
\]
Ici la constante d'intégration est le terme d'ordre 0 du développement limité, ie \(f(a)\).
\begin{center}
   \textit{
      Un développement limité de la dérivée nous donne donc \textbf{toujours}\+
      un développement limité de la fonction.
   }
\end{center}

\subsection*{\subsecstyle{Développements limités usuels {:}}}
\addcontentsline{toc}{subsection}{Développements limmités usuels}

Il est utile de connapitre les développements limités en \(0\) des fonctions usuelles:
\begin{center}
   \renewcommand{\arraystretch}{1.5}%
   \setlength\arrayrulewidth{0.8pt}

   \begin{tabular}{| c | c |}
   \hline
   \(e^x = 1 + x + \frac{x^2}{2} + \frac{x^3}{6} + \ldots + o(x^n)\) & \(\ln(1 + x) = x - \frac{x^2}{2} + \frac{x^3}{3} - \ldots + o(x^n) \)\\[1.5ex]
   \hline
   \((1+x)^\alpha = 1 + \binom{\alpha}{1} x + \binom{\alpha}{2}\frac{x^2}{2} + \binom{\alpha}{3}\frac{x^3}{6} + \ldots + o(x^n)\) & \(\frac{1}{1-x} = 1 + x + x^2 + x^3 + \ldots + o(x^n)\) \\[1.5ex]
   \hline
   \(\cos(x) = 1 - \frac{x^2}{2} + \frac{x^4}{24} - \ldots + o(x^n)\) & \(\sin(x) = x - \frac{x^3}{6} + \frac{x^5}{120} - \ldots + o(x^n)\) \\[1.5ex]   
   \hline
   \end{tabular}
\end{center} 

\chapter*{\chapterstyle{III --- Intégration}} % 99% Fini
\addcontentsline{toc}{section}{Intégration}

On définira dans ce chapitre l'intégration d'une fonction \textbf{bornée définie sur un segment}.\<

On appelle \textbf{subdivision} du segment \(\icc{a}{b}\) qu'on note généralement \(\sigma\) toute suite finie \((x_i)_{i \in \N}\) strictement croissante de premier terme \(x_0 = a\) et de dernier terme \(x_n = b\), on appelle \textbf{pas de la subdivision} la distance maximale entre deux éléments consécutifs de la subdivision, qu'on note généralement \(|\sigma|\).

\subsection*{\subsecstyle{Définition {:}}}
\addcontentsline{toc}{subsection}{Définition}

Soit \(f\) une fonction définie sur un segment \(\icc{a}{b}\), et \(\sigma = (x_i)_{i \in \N}\) une subdivision de ce segment, on note \(M_i := \sup\{f(x) \; ; \; x \in \icc{x_i}{x_{i+1}}\}\) le supremum de la fonction sur chaque intervalle (qui existe car on ne considère bien que les fonctions \textbf{bornées}).\<

Alors on appelle \textbf{somme de Darboux supérieure\footnote[1]{On définit de même \textbf{la somme de Darboux inférieure} associée à la subdivision qui est analogue à la différence que la hauteur des rectangle est donnée par l'infimum de la fonction sur la subdivision.} associée à la subdivision} la somme:
\[
   ^+S(f, \sigma) = \sum_{k=0}^{n} (x_{k+1} - x_k) \cdot M_k 
\] 
Cette somme correspond à \textbf{la somme de rectangles} chacun de base \(x_{k+1} - x_k\) et de hauteur le supremum de la fonction sur cette intervalle. Graphiquement, on a:
\begin{center}
   \begin{tikzpicture}[domain=0:2.55, xscale=1.9, yscale=0.6]
      \draw[color=DarkBlue1, thick] plot (\x,{((\x-1)^3 + 2)/1.5+1});
      \draw[-latex] (0,0) -- (3,0) node [right] {$x$};
      \draw[-latex] (0, 0) -- (0,5) node [above] {$y$};
      
      \draw[color = black] (0.35, 0.1) -- (0.35, -0.1) node[] at (0.35, -0.5) {$x_0$};
      \draw[dashed, color = black] (0.35, 0) -- (0.35, 2.3333) -- (1, 2.3333);
      \fill[black!30,nearly transparent] (0.35, 0) -- (0.35, 2.3333) -- (1, 2.3333) -- (1, 0) -- cycle;

      \draw[color = black] (1, 0.1) -- (1, -0.1) node[] at (1, -0.5) {$x_1$};
      \draw[dashed, color = black] (1, 0) -- (1, 2.417) -- (1.5, 2.417);
      \fill[black!30,nearly transparent] (1, 0) -- (1, 2.4166) -- (1.5, 2.4166) -- (1.5, 0) -- cycle;

      \draw[color = black] (1.5, 0.1) -- (1.5, -0.1) node[] at (1.5, -0.5) {$x_2$};
      \draw[dashed, color = black] (1.5, 0) -- (1.5, 3) -- (2, 3);

      \draw[dashed, color = black] (1.5, 3) -- (0, 3) node[color=black, left] {$\underset{\icc{x_2}{x_{3}}}{\sup}f$};
      \draw[color = black] (-0.05, 3) -- (0.05, 3);

      \fill[black!30,nearly transparent] (1.5, 0) -- (1.5, 3) -- (2, 3) -- (2, 0) -- cycle;

      \draw[color = black] (2, 0.1) -- (2, -0.1) node[] at (2, -0.5) {$x_3$};
      \draw[dashed, color = black] (2, 0) -- (2, 3.63) -- (2.25, 3.63) -- (2.25, 0);
      \fill[black!30,nearly transparent] (2, 0) -- (2,3.63) -- (2.25, 3.63) -- (2.25, 0) -- cycle;

      \draw[color = black] (2.25, 0.1) -- (2.25, -0.1) node[] at (2.25, -0.5) {$x_4$};

   \end{tikzpicture}
\end{center}
On appelle \textbf{l'intégrale supérieure}\footnote[2]{On définit de même \textbf{l'intégrale inférieure} qui est analogue à la différence que l'on considère l'aire \textbf{maximale} sur toutes les subdivisions possibles, ie on change l'infimum pour un supremum dans la définition.} de \(f\) la quantité:
\customBox{width=10cm}{
   \[^+S(f) = \inf\Bigl\{\,^+S(f, \sigma)\; ; \; \sigma \text{ est une subdivision de } \icc{a}{b}\Bigl\}\]
}
\begin{center}
   \textit{
      L'intégrale supérieure est la somme \textbf{minimale} obtenue en considérant \textbf{toutes les subdivisions possibles} du segment.
   }
\end{center}

Enfin on dit que \(f\) est \textbf{intégrable} si et seulement si \textbf{l'intégrale supérieure et inférieure} sont égales, et on la note:
\[
   \int_{a}^{b} f(t) \d t  
\]
L'ensemble des fonctions intégrables sur \(\icc{a}{b}\) est noté \(\mathscr{I}(\icc{a}{b})\).
\pagebreak

\subsection*{\subsecstyle{Caractérisation {:}}}
\addcontentsline{toc}{subsection}{Caractérisation}

Pour la suite, nous aurons besoin d'une propriété des intégrales supérieures et inférieures qui découle de leur nature de borne inférieure et supérieure\footnote[1]{Une somme supérieure pour une subdivision fixée est toujours plus grande que la borne inférieure de ces sommes pour \textbf{toutes} les subdivisions.}, en effet on a:
\customBox{width=7cm}{
   \({}^{-}S(f, \sigma) \leq {}^{-}S(f) \leq {}^{+}S(f) \leq {}^{+}S(f, \sigma)\)
}   
De cette propriété on peut déduire une caractérisation trés utile de l'intégrabilité\footnote[2]{En pratique, c'est celle ci qui permet de prouver qu'une classe de fonction est intégrable ou non.}, en effet une fonction \(f\) est intégrable si et seulement si pour tout \(\epsilon\) positif, \textbf{il existe une subdivision} \(\sigma\) telle que:
\customBox{width=5cm}{
   \({}^{+}S(f, \sigma) - {}^{-}S(f, \sigma) < \epsilon\)
} 
De cette proposition, nous pouvons alors montrer que plusieurs classes de fonctions sont intégrables, en particulier:
\begin{center}
   \textit{Toute fonction continue est intégrable.}\+
   \textit{Toute fonction monotone est intégrable. }
\end{center}

\subsection*{\subsecstyle{Propriétés {:}}}
\addcontentsline{toc}{subsection}{Propriétés}

Aprés avoir défini l'espace \(\mathscr{I}(\icc{a}{b})\) des fonctions intégrables, on peut considérer ses propriétés, en particulier on peut remarquer que cet espace est stable par somme et multiplication externe, en particulier:
\begin{center}
   \textit{L'espace des fonctions intégrables est un \textbf{sous-espace vectoriel} de l'espace des fonctions}
\end{center}
L'application d'intégration sur les fonctions intégrables est donc \textbf{une forme linéaire} sur cet espace, elle est \textbf{croissante} et vérifie \textbf{relation de Chasles}\footnote[3]{Pour \(c \in \icc{a}{b}\)} :
\[
   \int_{a}^{b} f(t) \d t = \int_{a}^{c} f(t) \d t + \int_{c}^{b} f(t) \d t
\]
  

L'intégrale vérifie aussi \textbf{l'inégalité triangulaire}, en effet on a:
\[ 
   \Bigl| \int_{a}^{b} f(t) \d t \Bigl| \leq \int_{a}^{b} |f(t)| \d t 
\]
  

\subsection*{\subsecstyle{Théorème fondamental du calcul intégral {:}}}
\addcontentsline{toc}{subsection}{Théorème fondamental du calcul intégral}

On peut maintenant énoncer le \textbf{théorème fondamental du calcul intégral}, on considère \(f\) une fonction continue sur un segment \(\icc{a}{b}\), et on définit la fonction \(F\) définie sur ce segment par:
\customBox{width=4cm}{
   \[F(x) = \int_{a}^{x} f(t) \d t\]
}   
Alors \(F\) est de classe\footnote[4]{Si \(f\) est seulement intégrable, alors \(F\) est seulement continue, c'est une version plus faible de ce théorème.} \(\mathscr{C}^1\) et est une \textbf{primitive} de \(f\), en particulier, c'est la seule qui s'annule en \(a\).\+
Par ailleurs, si \(f\) est intégrable et que l'on en connaît une primitive \(F\), alors on a:
\customBox{width=5cm}{
   \[\int_{a}^{b} f(t) \d t = F(b) - F(a)\]
}   

\underline{Exemple:} Calculons \(I = \int_{x}^{2x} f(t) \d t\) pour \(f = \cos(3t)\). On a \(I = F(2x) - F(x)\) pour \(F = \frac{\sin(3x)}{3}\) une primitive de \(f\) donc \(I = \frac{\sin(3x) - \sin(6x)}{3}\)

\subsection*{\subsecstyle{{Formule de la moyenne:}}}
\addcontentsline{toc}{subsection}{Formule de la moyenne}
Si \(f\) est continue sur le segment \(\icc{a}{b}\), alors il existe un \(c\) dans cet intervalle\footnote[1]{La preuve est immédiate, \(f\) est continue donc \(\int_{a}^{b} f(t) \d t = F(a) - F(b)\) et on conclut directement par le théorème des accroissements finis.} tel que:
\customBox{width=4.5cm}{
   \[f(c) = \frac{1}{b-a}\int_{a}^{b} f(t) \d t\]
}   
La quantité \(f(c)\) est appellée \textbf{valeur moyenne de la fonction sur le segment.}
Dans un cas plus général\footnote[2]{Même idée de démonstration, en utilisant la croissance de l'intégrale}, on peut même montrer que pour \(g\) continue et positive, alors on a l'existence d'un \(c\) tel que:
\[
   f(c)\int_{a}^{b} g(t) \d t = \int_{a}^{b} f(t)g(t) \d t
\]

\subsection*{\subsecstyle{Sommes de Riemann {:}}}
\addcontentsline{toc}{subsection}{Sommes de Riemann}

Soit \(f\) une fonction intégrable sur \(\icc{a}{b}\), \(\sigma_n = (x_k)_{k \leq n}\) une subdivision \textbf{régulière} et \(h = \frac{b - a}{n}\) \textbf{le pas de la subdivision}, alors on définit la \textbf{somme de Riemann} comme suit:
\customBox{width=7.5cm}{
   \[S_n = \frac{b - a}{n} \sum_{k=1}^{n}f\Bigl(a + k\frac{b - a}{n}\Bigl) = h \sum_{k=1}^{n}f(x_k)\]
}   

Alors la suite \((S_n)\) \textbf{converge} vers \(\int_{a}^{b} f(t) \d t\).\<

Une somme de Riemann représente la somme d'aires de rectangles dans le cas particulier d'une subdivision régulière et dans le cas où la hauteur des rectangle est simplement \(f(x_k)\). Cette propriété de convergence nous permet d'évaluer des séries gràce à des intégrales.\<

\underline{Exemple:} Dans le cas le plus courant où on est dans l'intervalle \(\icc{0}{1}\). Calculons la somme \(S := \sum_{k=0}^{+\infty} \frac{1}{n + k}\). On a:
\[
   S =  \frac{1}{n}\sum_{k=0}^{+\infty} \frac{1}{1 + \frac{k}{n}} = \frac{1}{n}\sum_{k=0}^{+\infty} f\Bigl(\frac{k}{n}\Bigl) = \int_{0}^{1} f(t) \d t = \ln(1 + x) \Bigl|_{0}^{1} = \ln(2)
\]

\subsection*{\subsecstyle{Techniques de calcul générales{:}}}
\addcontentsline{toc}{subsection}{Techniques de calcul générales}

Du théorème fondamental on peut déduire plusieurs formules utiles pour calculer des intégrales, dans toute cette section, \(f, g\) sont \textbf{continues} et \(\phi\) est \textbf{continue, dérivable à dérivée continue}\footnote[3]{De telle sorte que toutes les intégrales soient bien définies.}.\< 

En premier lieu on montre facilement la formule \textbf{d'intégration par parties} qui nous permet de caculer \textbf{l'intégrale d'un produit}\footnote[4]{On peut souvent trouver une relation de récurrence par IPP successives, ou alors annuler un facteur polynomial par les dérivations successives que permettent les IPP.} dans certains cas:
\customBox{width=8cm}{
   \[\int_{a}^{b} f(t)g'(t) \d y = f(t)g(t) \Bigl|_{a}^{b} \; - \int_{a}^{b} f'(t)g(t) \d t\]
}  

On montre par ailleurs la formule du \textbf{changement de variable}, elle nous permet de calculer \textbf{l'intégrale d'une composée}\footnote[5]{De gauche à droite, elle nous permet de poser \textbf{une nouvelle variable} comme fonction de la variable d'intégration, de droite à gauche on pose alors notre variable d'intégration comme \textbf{l'image d'une nouvelle variable par une fonction} (nécéssite alors la surjectivé sur le segment).} dans certains cas:
\customBox{width=6.5cm}{
   \[\int_{a}^{b} f(\varphi(t))\varphi'(t) \d t = \int_{\varphi(a)}^{\varphi(b)} f(x) \d x\]
}  


\subsection*{\subsecstyle{Intégration de fonctions complexes{:}}}
\addcontentsline{toc}{subsection}{Intégration de fonctions complexes}

Si on considère les fonctions trigonométriques et exponentielles, il est clair que l'on peut les voir comme la partie réelle ou imaginaire d'une fonction à valeur complexe.\+
Alors on peut utiliser les propriétés suivantes et profiter des propriétés de l'exponentielle:
\customBox{width=12.5cm}{
   \[   \int_{a}^{b} \mathscr{R}(f(t)) \d t = \mathscr{R}\Bigl(\int_{a}^{b} f(t) \d t\Bigl) 
   \quad\quad\quad\quad
   \int_{a}^{b} \mathscr{I}(f(t)) \d t = \mathscr{I}\Bigl(\int_{a}^{b} f(t) \d t\Bigl) \]
}  

\underline{Exemple:} Calculons \(I := \int e^{-x}\cos(x) \d x\)
\[
   I = \int e^{-x}\mathscr{R}(e^{ix}) \d x = \int \mathscr{R}(e^{x(i- 1)}) \d x = \mathscr{R}(\int e^{x(i - 1)} \d x) = \mathscr{R}(\frac{1}{i - 1}e^{x(i - 1)}) = \frac{e^{-x}}{2}(\sin(x) - \cos(x))
\]
On remarque ici que l'intégrale \(\int e^{\alpha x} \d x\) se comporte de manière identique si \(\alpha\) est réel ou complexe.

\subsection*{\subsecstyle{Décomposition en éléments simples{:}}}
\addcontentsline{toc}{subsection}{Décomposition en éléments simples}

Si on considère les primitives de fractions rationnelles de la forme \(\frac{P}{Q}\), on peut montrer qu'elles peuvent se décomposer en une somme \textbf{d'éléments simples} particulièrement faciles à intégrer. Cette partie se consacre à la décomposition d'une fraction \(\frac{P}{Q}\) en éléments simples\footnote[1]{On considère \(d^{\circ}(P) < d^{\circ}(Q)\), dans le cas opposé, on peut s'y ramener par division euclidienne.}.
\begin{center}
   \begin{itemize}
      \item La première étape de la décomposition consiste à \textbf{factoriser} \(Q\) en polynômes irréductibles.
      \item Tout facteur \((aX + b)^n\) produira alors une somme de \(n\) éléments simples de première espèce:
      \[
         \frac{\lambda_1}{aX+b} + \frac{\lambda_2}{(aX+b)^2} + \ldots + \frac{\lambda_n}{(aX+b)^n}
      \] 
      \item Tout facteur \((aX^2 + bX + c)^n\) produira alors une somme de \(n\) éléments simples de seconde espèce:
      \[
         \frac{\lambda_1X+\beta_1}{aX^2+bX+c} + \frac{\lambda_2X+\beta_2}{(aX^2+bX+c)^2} + \ldots + \frac{\lambda_nX+\beta_n}{(aX^2+bX+c)^n}
      \]
   \end{itemize}
\end{center}

Finalement, on obtient une somme d'éléments simples avec des coefficients indéterminés aux numérateurs. Pour déterminer ces coefficients, plusieurs méthodes sont possibles, réduire au même dénominateur et identifier, évaluer en des valeurs annulatrices ou encore utiliser \textbf{la méthode du cache}\footnote[2]{Cette méthode permet de calculer les coefficients d'un élément dont le dénominateur est de plus haut degré de son groupe, elle consiste à multiplier l'égalité obtenue par le dénominateur en question, puis à évaluer en une de ces racines, annulant ainsi \textbf{tout les éléments simples} sauf celui dont on a choisi le dénominateur, cette procédure permet alors de trouver les coefficients de cet élémént.}.\< 

\underline{Exemple:} Décomposons \(\frac{2X+1}{(X-1)(X+1)^2(X^2+1)}\)\+
On a:
   \begin{align*}
      \frac{2X+1}{(X-1)(X+1)^2(X^2+1)} &= \frac{\lambda_1}{X-1} + \frac{\lambda_2}{X+1}+ \frac{\lambda_3}{(X+1)^2} + \frac{\lambda_4X+ \lambda_5}{X^2+1}
   \end{align*}
Utilisons la méthode du cache sur l'élement en \(X-1\), on multiplie par \(X-1\) de chaque coté et on évalue en la racine \(1\) ce qui nous donne \(\lambda_1 = \frac{3}{8}\).\+
Utilisons la méthode du cache sur l'élement en \((X+1)^2\), on multiplie par \((X+1)^2\) de chaque coté et on évalue en la racine \(-1\) ce qui nous donne \(\lambda_3 = \frac{1}{4}\).\+
Utilisons la méthode du cache sur l'élement en \(X^2+1\), on multiplie par \(X^2+1\) de chaque coté et on évalue en la racine \(i\) ce qui nous donne \(\lambda_4i+\lambda_5 = \frac{-1}{4}i - \frac{3}{4}\) donc \(\lambda_4=\frac{-1}{4}\) et \(\lambda_5=\frac{-3}{4}\).\<

Pour finir pour trouver le dernier coefficient, on peut par exemple réduire au même dénominateur et identifier, pour conclure on a:
\begin{align*}
   \frac{2X+1}{(X-1)(X+1)^2(X^2+1)} &= \color{BrightRed1}\underbrace{\color{black}\frac{3}{8(X-1)} + \frac{-1}{8(X+1)}+ \frac{1}{4(X+1)^2}}_\text{Première espèce} \color{black}+\color{BrightRed1}\underbrace{\color{black} \frac{-X - 3}{4(X^2+1)}}_\text{Seconde espèce} 
\end{align*}

\subsection*{\subsecstyle{Intégration des éléments simples{:}}}
\addcontentsline{toc}{subsection}{Intégration des éléments simples}

Les éléments simples \textbf{de première espèce} sont ceux de la forme ci-dessous, ie dont le dénominateur est \textbf{une puissance de polynôme irreductible de premier degré}, pour les intégrer, un changement de variable affine conclut:
\customBox{width=3.5cm}{
   \[\int_{a}^{b} \frac{1}{(at+b)^n}  \d t \]
}  

Les éléments simples \textbf{de seconde espèce}, sont ceux dont le dénominateur est \textbf{une puissance de polynôme irreductible de second degré}, pour les intégrer, on se ramène à la forme ci-dessous, et un calcul par récurrence ou un changement de variable trigonométrique conclut\footnote[1]{On peut faire le changement de variable \(u = \tan(x)\) ou alors trouver une relation de récurrence en faisant une IPP à partir de \(I_{n+1}\).}:
\customBox{width=3.5cm}{
   \[\int_{a}^{b} \frac{1}{(1+t^2)^n}  \d t \]
}  

Pour se ramener à la forme ci-dessus, pour \(E = \frac{P}{Q^n}\) un élément de seconde espèce, on suit la méthode suivante:\+
\begin{description}
   \item[$\bullet$] On fait \textbf{la division euclidienne de \(P\) par \(Q'\)\footnote[2]{\(Q'\) sera nécessairement de degré 1.}}.
   \item[$\bullet$] On obtient alors une intégrale simple et un reste de la forme \(\frac{1}{Q^n}\).
   \item[$\bullet$] On ramene le reste à la forme souhaitée en mettant \(Q\) sous \textbf{forme canonique}.
\end{description}
\underline{Exemple:} L'élément de seconde espèce de l'exemple précédent est \(E = \frac{-X - 3}{4(X^2+1)}\), on fait la division euclidienne du numérateur par la dérivée du dénominateur \(8X + 4\) et on a:
\[
   E = \frac{\frac{-1}{8}(8X - 4) + \frac{-5}{2}}{4(X^2+1)} =  \frac{-1}{8} \cdot \frac{2X - 1}{X^2+1} - \frac{-5}{8(X^2 +1)}
\]
Puis on peut intégrer:
\[
   \int E(x) \d x = \frac{-1}{8}  \color{DarkBlue1}\underbrace{\color{black} \int \frac{2x - 1}{x^2+1} \d x}_\text{Facile}\color{black}  - \frac{-5}{8} \color{DarkBlue1}\underbrace{\color{black}\int \frac{1}{(x^2 +1)} \d x}_\text{Forme souhaitée}
\]

\subsection*{\subsecstyle{Parité{:}}}
\addcontentsline{toc}{subsection}{Parité}

Soit \(a \in \R^+\), il est possible d'exploiter la parité de la fonction\footnote[3]{Par ailleurs, on remarque qu'il est possible de décomposer une fonction en somme d'une fonction paire et d'une fonction impaire, en effet on a:
\(f(x) = \frac{f(x) + f(-x)}{2} + \frac{f(x) - f(-x)}{2}\)} à intégrer pour faciliter les calculs, en effet on a:
\begin{itemize}
   \item Si la fonction est \textbf{paire}, on a:
   \customBox{width=5cm}{
      \[\int_{-a}^{a} f(t) \d t = 2\int_{0}^{a} f(t) \d t\]
   }  
   \item Si la fonction est \textbf{impaire}, on a:
   \customBox{width=3.5cm}{
      \[\int_{-a}^{a} f(t) \d t = 0\]
   }  
\end{itemize}
